\subsection{CP-02: PSE de Wompi, redirección y consulta de estado}
\label{sec:cp02}
\label{anexo:cp02}

\FichaCaso
  {CP-02}
  {Equipo Kit Pagos Colombia}
  {09/11/2026}
  {PSE de Wompi: redirección y consulta de estado}
  {Pendiente (\#174)}
  {Pendiente (\#174)}
  {Crear un pago con PSE en Wompi y comprobar que el SDK distingue una redirección requerida de una transacción resuelta: la creación devuelve la URL del banco y el identificador de la pasarela con estado nativo \texttt{PENDING}; la consulta posterior devuelve \texttt{APPROVED}. El caso verifica además que el código de banco se obtiene con \texttt{getPseBanks()} y no está escrito a mano.}
  {Funcional (sistema, caja negra).}
  {\textbullet\ SDK compilado y dependencias de \texttt{examples/} instaladas.\newline
   \textbullet\ API de Simulación arriba en \texttt{http://localhost:3000}.\newline
   \textbullet\ El ejemplo pide la lista de bancos con \texttt{getPseBanks()} y usa la primera entidad; no fija un código.\newline
   \textbullet\ Documento del pagador disponible, porque PSE lo exige.}
  {Local, sin credenciales reales. Simulador local en el puerto 3000; SDK 0.7.0-rc.2; monto \texttt{150000.00 COP}; banco elegido de la lista del simulador (primera entidad, código \texttt{1}, ``Banco que aprueba''). Node.js 20.}

\subsubsection{Objetivo}

Comprobar que el resultado de la creación conserva la diferencia entre \texttt{TRANSACTION} y \texttt{REDIRECT\_REQUIRED}, que la redirección trae una URL de banco y el identificador nativo, y que la consulta por identificador resuelve el pago sin que el comercio conozca el vocabulario de Wompi.

\subsubsection{Procedimiento}

\begin{TablaPasos}{Pasos del caso CP-02}
  1 &
  Verificar que el simulador responde. &
  \texttt{cd simulator-api \&\& npm run dev}; luego \texttt{GET}\ \url{http://localhost:3000/health}. &
  Respuesta 200 con estado \texttt{ok}. &
  & & \\ \hline
  2 &
  Ejecutar el ejemplo de PSE de Wompi. &
  \texttt{cd examples \&\& npm run simulate:wompi-pse}. &
  La salida imprime la solicitud con el documento del pagador y el código de banco elegido; luego muestra una redirección pendiente con la URL del banco, el identificador de la pasarela y estado nativo \texttt{PENDING}. No imprime una transacción resuelta en este paso. &
  & & \\ \hline
  3 &
  Verificar el discriminante del resultado. &
  Inspección de la salida y del código de \texttt{simulate-wompi-pse.ts}. &
  El ejemplo entra por la rama \texttt{result.outcome === "REDIRECT\_REQUIRED"} y lee \texttt{result.redirect}; el compilador no permite tratarlo como una transacción sin comprobar la rama. &
  & & \\ \hline
  4 &
  Consultar el estado después de la redirección. &
  El identificador impreso en el paso 2, con \url{kitPagos.getPaymentStatus(id)}. &
  La consulta devuelve el mismo identificador y estado normalizado \texttt{APPROVED}; el estado nativo es \texttt{APPROVED}. &
  & & \\ \hline
  5 &
  Verificar el origen del código de banco. &
  Revisión de \texttt{simulate-wompi-pse.ts}. &
  El código sale de \texttt{kitPagos.getPseBanks()} y se usa el primero de la lista; no hay un código numérico escrito a mano. &
  & & \\ \hline
  6 &
  Repetir la corrida. &
  El mismo comando del paso 2. &
  La segunda corrida reproduce la redirección, el \texttt{PENDING} y el \texttt{APPROVED}; el proceso termina con código 0. &
  & & \\ \hline
\end{xltabular}
\end{TablaPasos}

\subsubsection{Criterio de aceptación}

\begin{enumerate}
  \item La creación devuelve \texttt{outcome = REDIRECT\_REQUIRED}, no una transacción resuelta.
  \item La redirección incluye una URL de banco no vacía y un \texttt{gatewayTransactionId} consultable.
  \item El estado nativo de la creación es \texttt{PENDING} y el de la consulta es \texttt{APPROVED}.
  \item El código de banco se obtuvo de \texttt{getPseBanks()} y corresponde a una entidad de la lista del simulador.
  \item El ejemplo termina con código de salida 0 en las dos corridas.
\end{enumerate}

\subsubsection{Observaciones de diseño}

Contra el sandbox real de Wompi este flujo no se puede ejercitar: el sandbox publica la URL de redirección en el mismo instante en que resuelve el pago, así que cuando la URL existe ya no sirve y no hay ventana para redirigir. La API de Simulación reproduce el orden real —primero la redirección, después la consulta— y por eso es el único ambiente donde el caso tiene sentido. Esta limitación está documentada en el punto 43 del \texttt{architecture-log.md} y es una razón técnica para ejecutar el caso contra el simulador.

\subsubsection{Trazabilidad}

RF-01, RF-03, RF-06 y RF-19; HU-01, HU-02 y HU-12.
